\section{The endpoint gain and the closed surface}
\label{sec:final-proof}

The branch of Theorem~\ref{thm:weak-branch} was constructed with the
weight $H^\beta$, because the reference inverse is available only for
exponents below one. This already suffices for the geometry: every
element of $\mathcal X_\beta$ is flat at the endpoints. In this
section we first show that the corrections to the metric are in fact
of size $q^2H^{2\beta}$, which quantifies how the final metric
differs from the reference one, and then we state the annular theorem
in its precise form and prove Theorems~\ref{thm:annulus-intro}
and~\ref{thm:main}.

The gain comes from two identities for angular means. The first
recovers the mean of the height from an expression which is quadratic
in the oscillating part. The second recovers the controls $\kappa$
from that mean and from another quadratic expression. A product of
two factors of size $H^\beta$ has size $H^{2\beta}$, and
$2\beta>1$. No new iteration in a stronger space is needed. In this
section a quantity of size $H^\beta$ is called \emph{weak} and one of
size $H^{2\beta}$ \emph{strong}; both are flat at the endpoints.

\subsection{Two identities for means}

\paragraph{The mean of the height.}
Let $\mathfrak F_w(U)$ be the connection period
\eqref{eq:cart-mean-path-flux}. By
\eqref{eq:mean-period-derivative} its derivative at $U=0$ is
$w\av{U_r}$. Let
\begin{equation}\label{eq:posterior-period-bilinear}
 \mathcal B_{w,U}(v_1,v_2)
 =\int_0^1(1-t)\,D_U^2\mathfrak F_w(tU)[v_1,v_2]\,dt
\end{equation}
be its Taylor remainder of second order. For a solution the period
vanishes identically, because its radial derivative is the mean of the
residual and it is zero at the endpoints. Hence, with
$u^{\rm av}=\av u$,
\begin{equation}\label{eq:mean-recovered-derivative}
 0=\mathfrak F_w(qu)=q\,w\,(u^{\rm av})'
   +q^2\,\mathcal B_{w,qu}(u,u),
 \qquad\text{that is,}\qquad
 (u^{\rm av})'=-\frac qw\,\mathcal B_{w,qu}(u,u) .
\end{equation}

\paragraph{The mean of the return residual.}
Use the additive coordinates $\mathbf v=(Z,\kappa,a)$ with $Z=qu$
and $h_L=\id+a$, and put
\[
 \mathcal C(\mathbf v)=\av{\mathfrak e(\mathbf v)},
 \qquad
 \mathcal B^\circ_{\mathbf v}(v_1,v_2)
 =\int_0^1(1-t)\,D^2\mathcal C(t\mathbf v)[v_1,v_2]\,dt .
\]
By Section~\ref{sec:solving-the-equation} the derivative of
$\mathcal C$ at $\mathbf v=0$ is the identity on $\kappa$, zero on
the odd increments of the charts, and $R_0^\circ$ on the mean of the
height. At a solution therefore
\begin{equation}\label{eq:weak-posterior-control}
 0=\kappa+q\,R_0^\circ u^{\rm av}
       +\mathcal B^\circ_{\mathbf v}(\mathbf v,\mathbf v) .
\end{equation}

Formulas \eqref{eq:mean-recovered-derivative} and
\eqref{eq:weak-posterior-control} already explain the exponent
$2\beta$. What remains is to verify the corresponding estimates in
the stronger topology, with all derivatives in the parameter, so that
the conclusion holds for the whole branch.

\subsection{Two weak factors give one strong remainder}

Fix a convex neighbourhood of $(U,\nu,\kappa,a)=(0,0,0,0)$ in the
spaces of weight $\beta$, small in one norm of finite order, on which
$D$ and $Q/w$ are positive and the charts are uniformly increasing.
Both Taylor remainders above are integrals over segments in this
neighbourhood; the map $\mathcal C$ does not depend on $\nu$.

\begin{lemma}[Quadratic remainders in the doubled weight]
\label{lem:posterior-two-weight}
The maps $(w,U,v_1,v_2)\mapsto w^{-1}\mathcal B_{w,U}(v_1,v_2)$ and
$(\mathbf v,v_1,v_2)\mapsto\mathcal B^\circ_{\mathbf v}(v_1,v_2)$
are smooth tame with values in the space $\mathcal X^0_{2\beta}$ of
radial functions and in $\mathcal K_{2\beta}$, respectively. For every
number $k$ of derivatives with respect to the state there is a loss
$d_k$, independent of the order $s$, such that
\begin{equation}\label{eq:posterior-multilinear-bound}
 \begin{split}
 &\bigl\|D^k\mathcal B[\eta_1,\ldots,\eta_k]\,
                         (\eta_{k+1},\eta_{k+2})\bigr\|_{2\beta,s}\\
 &\quad\le C_{s,k}\Bigl[\,
   \sum_{i=1}^{k+2}\|\eta_i\|_{s+d_k}
                    \prod_{j\ne i}\|\eta_j\|_{b+d_k}
   +(1+\|\mathbf y\|_{s+d_k})
                    \prod_j\|\eta_j\|_{b+d_k}\Bigr],
 \end{split}
\end{equation}
where $\mathcal B$ is either of the two maps, $\mathbf y$ is the
state, the norms of the arguments are those of weight $\beta$, and
compact scalars are measured by their absolute values. The same bound
holds for all derivatives of $\mathcal C$ of order at least two.
\end{lemma}

\begin{proof}
\step{1: counting weak factors}
On the component $m$, every increment of a control or of a chart, and
every derivative of positive order of a height, carries a factor
$H_m^\beta$. The endpoint constants of heights are harmless: only
derivatives of positive order enter, and those of $\chi_L$ are
supported in a compact positive region and vanish on the supports of
all controls on negative components. Differentiation does not
destroy a weak factor: for $u=H^\beta\widetilde u$,
\[
 H^{-\beta}\,\Dr u=\Dr\widetilde u+\beta\,(\Dr\log H)\,\widetilde u ,
\]
and $\Dr\log H$ is bounded. Products and restrictions to the pieces
cost fixed numbers of derivatives and reference factors
$\mathcal E_m^{C_{s,k}}$. So it suffices to show that every term of
the expressions in question contains two weak factors, one from each
of the two distinguished arguments: these give the weight
$H_m^{2\beta}$ of the output; the highest derivative is placed on one
argument or on one coefficient of the state; the complete product of
coefficients is absorbed into one power of $M$; and finite overlap
together with $\|ab\|_{\ell^2}\le\|a\|_{\ell^2}\|b\|_{\ell^\infty}$
sums over the components.

\step{2: the period}
With $k=w+U_r$, $p=U_\theta/r$ and $E=1+w^2$, the integrand of
$\mathfrak F_w(U)$ is
\[
 \frac{1+kp_\theta/(1-p^2)}{\sqrt{E(1-p^2)-k^2}}-1 ,
\]
a smooth function of $w$ and of the derivatives
$(U_r,U_\theta/r,U_{\theta\theta}/r)$ on the neighbourhood. Each
second derivative with respect to $U$ contains one such derivative
from each argument, and further differentiation with respect to the
state does not remove them. The factor
$w^{-1}=w_*^{-1}\exp(-\nu\chi_*-\sum_m\kappa_m\chi_m)$ is part of the
product of coefficients; it costs a coefficient power, not an
endpoint exponent.

\step{3: flows}
For the map $\mathcal C$ consider a normalized characteristic flow
$\dot\Theta=G(t,\Theta;\mathbf y)$, in the form
\eqref{eq:fixed-strip-root-flow} near a root. Its first and second
variations with respect to the state satisfy
\[
 \dot\Theta_i=G_\theta\,\Theta_i+G_i,
 \qquad
 \dot\Theta_{12}=G_\theta\,\Theta_{12}
   +G_{\theta\theta}\,\Theta_1\Theta_2
   +(G_1)_\theta\,\Theta_2+(G_2)_\theta\,\Theta_1+G_{12} .
\]
The integrating factor of the homogeneous equation is bounded above
and below, as in the proof of
Proposition~\ref{prop:uniform-physical-geometry}. The source $G_i$
contains one weak factor and $G_{12}$ two, so by variation of
constants $\Theta_i$ contains one and $\Theta_{12}$ two. Composition
of flows does not change the weight of the component.

\step{4: inverses}
If $\iota=h^{-1}$, then
\[
 \iota_i=-\frac{h_i\circ\iota}{h'\circ\iota},\qquad
 \iota_{12}=-\frac{h_{12}\circ\iota
 +(h_1'\circ\iota)\,\iota_2+(h_2'\circ\iota)\,\iota_1
 +(h''\circ\iota)\,\iota_1\iota_2}{h'\circ\iota} .
\]
The denominator is bounded below, and every term of the second
numerator contains both variations. Apply this to $f^{-1}$ and
$h_R^{-1}$ in \eqref{eq:full-height-one-chart}; taking even or odd
parts and means preserves the weights.

\step{5: higher variations}
Label the directions of differentiation by a finite set $I$. In the
chain rule for $G(t,\Theta;\mathbf y)$, a subset $J\subset I$ acts
directly on the state and the blocks $B$ of a partition of
$I\setminus J$ act on the flow, giving terms
\[
 \bigl(D_{\mathbf y}^{|J|}\partial_\theta^{\ell}G\,
                   [\eta_j:j\in J]\bigr)\circ\Theta\;
                            \prod_{B}\Theta_B ,
\]
where $\ell$ is the number of blocks. Keep the term
$G_\theta\,\Theta_I$ on the left side; all other terms involve only
variations of lower order and contain every label exactly once. By
induction both distinguished factors are retained, and the identity
for inverses has the same structure. The additional weak factors are
bounded in low norms. The tame estimates for products and
compositions put only one factor at the high order, and the formulas
for the coefficients, the regularization at the roots and the
restrictions have fixed order; $k$ variations add a cost depending on
$k$ but not on $s$. This proves
\eqref{eq:posterior-multilinear-bound}.

\step{6: smoothness}
These are estimates for the actual derivatives. Taylor's formula with
integral remainder and one more variation gives continuity and
differentiability in the norms of weight $2\beta$, and by iteration
smoothness in the stronger Fr\'echet topology, not only pointwise
decay.
\end{proof}

\subsection{Recovering the mean and the controls}

\begin{lemma}[A primitive keeps the endpoint exponent]
\label{lem:weighted-primitive}
Let $\gamma>0$. For a radial function $f\in\mathcal X^0_\gamma$ let
$\mathscr Pf$ be its primitive which vanishes at the right endpoint.
Then $\mathscr Pf=v+a_L\chi_L$ with $v\in\mathcal X^0_\gamma$ and
$a_L\in\R$, and $\mathscr P:\mathcal X^0_\gamma\to\mathcal X_\gamma$
is a tame linear map on radial functions.
\end{lemma}

\begin{proof}
Near an endpoint, orient $\varrho$ towards it. Then
$(\log H)'\to-1$ in the variable $\varrho$, so far enough out
\[
 \frac{d}{d\varrho}\log H^\gamma\le-\frac\gamma2 :
\]
the ratio of $H^\gamma$ at $\sigma\ge\varrho$ to its value at
$\varrho$ is at most $e^{-\gamma(\sigma-\varrho)/2}$. Near the right
endpoint $\mathscr Pf(\varrho)=-\int_\varrho^\infty f\,\ell\,d\sigma$,
where the factor $\ell\le d_\partial$ comes from $dr=\ell\,d\varrho$
and is bounded. So the kernel of $\mathscr P$, conjugated by the
weight $H^{-\gamma}$, is bounded by $Ce^{-\gamma(\sigma-\varrho)/2}$
for $\sigma\ge\varrho$ and vanishes otherwise. Its integrals in
either variable are bounded, and Schur's test gives the bound in
$L^2(d\varrho)$. One derivative of the primitive is $\ell f$; further
derivatives with respect to $\varrho$ fall on $f$, on the weight and
on $\ell$, whose logarithmic derivatives are bounded. A factor $M^c$
with $c\ge0$ causes no loss, since $M(\varrho)\le M(\sigma)$ for
$\sigma\ge\varrho$. Near the left endpoint the same
argument applies to $\mathscr Pf-a_L$, where
$a_L=-\int_{r_-}^{r_+}f\,dr$ is a continuous functional of $f$,
because this difference is the primitive which vanishes at the left
endpoint. On the compact middle part integration is bounded in every
Sobolev norm, and the derivative of $\chi_L$ is supported there.
\end{proof}

\paragraph{The mean.}
By \eqref{eq:mean-recovered-derivative} the mean of the height on the
solution branch is given by the explicit map
\begin{equation}\label{eq:rewrite-mean-recovery}
 \mathcal U(\bx,q)
 =-q\,\mathscr P\bigl[w^{-1}\mathcal B_{w,qu}(u,u)\bigr],
 \qquad u_q^{\rm av}=\mathcal U(\bx(q),q) .
\end{equation}
By Lemma~\ref{lem:posterior-two-weight}, $\mathcal U$ is smooth with
values in $\mathcal X_{2\beta}$, and it is $O(q)$ in every norm; in
particular the difference of the endpoint values of $u_q$ is $O(q)$.

\paragraph{The controls.}
The operator $R_0^\circ$ of
\eqref{eq:rewrite-reference-circle-row} involves only the derivative
of its argument. It maps the part of $\mathcal U$ in
$\mathcal X^0_{2\beta}$ to $\mathcal K_{2\beta}$ with a fixed loss,
and annihilates the multiple of $\chi_L$, whose derivative is
supported in a positive component. By
\eqref{eq:weak-posterior-control},
\begin{equation}\label{eq:rewrite-control-recovery}
 \kappa_q=-q\,R_0^\circ\,\mathcal U(\bx(q),q)
            -\mathcal B^\circ_{\mathbf v(q)}(\mathbf v(q),\mathbf v(q)) .
\end{equation}
Since $\mathbf v(q)=O(q)$ on the branch, both terms on the right are
$O(q^2)$ in every norm of $\mathcal K_{2\beta}$. The identity also
shows that $q\mapsto\kappa_q$ is smooth with values in
$\mathcal K_{2\beta}$, and that $\dot\kappa_0=0$, where the dot
denotes the derivative with respect to $q$.

\paragraph{The compact parameter.}
Differentiating the equation $\mathcal F_q(w_q,u_q)=0$ at $q=0$,
\[
 \mathscr L_{w_*}\dot u_0+\dot\nu_0\,k_*+K_*\,\dot\kappa_0
       +\partial_q\mathcal F_q(w_*,\Phi_*)\big|_{q=0}=0,
 \qquad \ell_*\dot u_0=0 .
\]
The last term is quadratic in the derivatives of $\Phi_*$, which
contain only the angular frequency two; products of $\cos2\theta$ and
$\sin2\theta$ contain only the frequencies zero and four. So this
term has no component in the second mode. Since also
$\dot\kappa_0=0$, the second mode of the equation is the homogeneous
bordered equation of Proposition~\ref{prop:rad-fixed-theorem}, whose
only solution is zero. Hence $\dot\nu_0=0$ and $\nu_q=O(q^2)$.

\begin{figure}[htbp]
\centering
\begin{tikzpicture}[x=6.6cm,y=1cm,font=\small,>=Stealth]
\draw[->] (0,0)--(1.75,0) node[right] {endpoint exponent};
\fill[blue!8] (.5,-.12) rectangle (1,.12);
\foreach \x/\lab in {0/0,.5/{1/2},.75/{\beta=3/4},1/1,1.5/{2\beta=3/2}}
 {\draw (\x,-.12)--(\x,.12); \node[below] at (\x,-.16) {$\lab$};}
\draw[->,thick,orange!85!black] (.75,.4) to[bend left=30] (1.5,.4);
\node[above] at (1.12,.95) {two weak factors in each remainder of a mean};
\node[align=center] at (.75,-1.05) {admissible range\\$1/2<\beta<1$};
\end{tikzpicture}
\caption{The two roles of the exponent $\beta=\tfrac34$. Exponents
near $\beta$ and below one are needed for the reference inverse. The
equations for the means multiply two factors of weight $\beta$ and
give an exponent above one, which yields the bound for the metric in
Theorem~\ref{thm:new-annular-height}. Flatness of the solution needs
only a positive exponent.}
\label{fig:exponent-gain}
\end{figure}

\subsection{The annular theorem}

For a profile $w$ let $h_w$ denote the radial function with
$h_w'=w$, with a fixed normalization of the additive constant.

\begin{theorem}[A nonradial height with flat ends]
\label{thm:new-annular-height}
Let $w_*$, $\phi_*$, $\chi_*$ be as in
Proposition~\ref{prop:rad-fixed-theorem} and let $\chi_m$ be the
controls of Lemma~\ref{lem:quantitative-optical-controls}. There are
$q_0>0$ and a smooth family
\[
 (u_q,\nu_q,\kappa_q,h_{L,q}),\qquad |q|<q_0,
\]
through $(\Phi_*,0,0,\id)$ at $q=0$, such that the profile and the
height
\[
 w_q=w_*\exp\Bigl(\nu_q\chi_*+\sum_m\kappa_{q,m}\chi_m\Bigr),
 \qquad z_q=h_{w_q}+q\,u_q
\]
have the following properties, where
$g_q=(1+w_q^2)\,dr^2+r^2\,d\theta^2$.
\begin{enumerate}
\item \textup{(Equation.)} $|dz_q|_{g_q}<1$ and
 $g_q-dz_q\otimes dz_q$ is flat; equivalently
 $\mathcal N_{h_{w_q}}(z_q)=0$ with $D>0$.
\item \textup{(Profile.)} $w_q>0$ on $(r_-,r_+)$, and $w_q$ extends
 smoothly by zero across both endpoints. The zeros of $w_q'$ are
 those of $w_*'$, and $w_q=w_*$ near each of them.
\item \textup{(Height.)} $z_q$ is even and $\pi$-periodic in
 $\theta$, and has a constant extension at each endpoint; the two
 constants may differ. Moreover $\ell_*u_q=1$, so $z_q$ is not
 radial for $q\ne0$.
\item \textup{(Returns.)} On every negative component $m$, the two
 characteristic return maps $f_{q,m,\pm}$ of the linearized equation
 at $z_q$ satisfy
 \begin{equation}\label{eq:annular-return-condition}
  h_{R,q,m}\circ f_{q,m,\pm}\circ h_{L,q,m}^{-1}=R_{\pm T_m}
 \end{equation}
 for odd increasing charts commuting with the translation by $\pi$,
 where $T_m$ is the travel angle of $w_*$.
\item \textup{(Estimates.)} For every $s$,
 \begin{equation}\label{eq:annular-branch-estimates}
  \|u_q-\Phi_*\|_{\mathcal X_\beta,s}\le C_s\,|q|,\qquad
  |\nu_q|+\|\kappa_q\|_{\mathcal K_{2\beta},s}\le C_s\,q^2 ,
 \end{equation}
 and the family is smooth in $q$ for these topologies. In particular,
 for all $a\ge0$ and $\epsilon>0$,
 \[
  |\partial_r^a(w_q-w_*)|\le C_{a,\epsilon}\,q^2\,H^{1-\epsilon}
 \]
 near the endpoints, and $w_q-w_*=O(q^2)$ in every $C^a$ norm.
\end{enumerate}
\end{theorem}

\begin{proof}
The family is the branch of Theorem~\ref{thm:weak-branch}. Choose
$q_0$ so that $[-q_0,q_0]$ lies in its domain; every seminorm of the
branch is bounded on this compact interval, with constants depending
on the order but with $q_0$ independent of the order.

\emph{Estimates.} The first bound in
\eqref{eq:annular-branch-estimates} is contained in
Theorem~\ref{thm:weak-branch}; the second and the smoothness in the
stronger topology follow from
\eqref{eq:rewrite-mean-recovery}, \eqref{eq:rewrite-control-recovery}
and the bound for $\nu_q$. On the tails $\chi_*$ vanishes, and
differentiating the formula for $w_q$ gives
$|\partial_r^a(w_q-w_*)|\le C_a\,q^2\,H^{2\beta}E_0^{C_a}$; since
$2\beta>1$, \eqref{eq:full-material-envelope} gives the stated bound.
On the compact middle part use the smoothness of the family and
$\nu_q=O(q^2)$.

\emph{Profile.} The controls vanish near every turning circle, so the
germ of the metric there is unchanged, and by
Proposition~\ref{prop:uniform-physical-geometry}(i) the sign of
$w_q'$ is that of $w_*'$ elsewhere. Flatness of $w_q$ follows from
that of $w_*$ and of the controls.

\emph{Equation and returns.} At a zero of $\mathbf R$ both residual
corrections vanish, so $\widehat L$ is the true linearization and its
exact returns are \eqref{eq:annular-return-condition}. The inequality
$D>0$ of Proposition~\ref{prop:uniform-physical-geometry}(i) is
equivalent to $|dz_q|_{g_q}<1$, and $\mathcal A=0$ means that the
curvature residual vanishes, which by
Proposition~\ref{prop:scalar-core} is the flatness of the remainder
metric.

\emph{Height.} The symmetry is that of the function spaces. Since
$\ell_*u_q=1$, the second angular mode of $u_q$ is not zero. By
Lemma~\ref{prop:full-physical-tame} every element of
$\mathcal X_\beta$ has constant extensions at the endpoints; the same
holds for all derivatives in $q$. Finally $h_{w_q}$ has constant
extensions because $w_q$ is flat.
\end{proof}

\subsection{Proof of Theorems~\ref{thm:annulus-intro} and~\ref{thm:main}}

The construction involves two parameters with different roles, and it
is worth separating them once more. The small parameter $q$ was used
to construct one nonradial isometric realization of one rotationally
symmetric metric. The phase $\tau$ moves that realization.

\begin{proof}[Proof of Theorem~\ref{thm:annulus-intro}]
Fix $q$ with $0<|q|<q_0$ and let $w=w_q$, $z=z_q$ be given by
Theorem~\ref{thm:new-annular-height}. Part~(i) of that theorem is
condition~(i); parts~(ii) and~(iii) give conditions~(ii) and~(iii).
\end{proof}

\begin{proof}[Proof of Theorem~\ref{thm:main}]
By Section~\ref{sec:geometric-closure},
Theorem~\ref{thm:annulus-intro} implies Theorem~\ref{thm:main}. In
detail: extend $w$ by zero and $z$ by its limiting constants to a
slightly larger annulus. The remainder metric is Euclidean on the
added collars, so the extended data satisfy the hypotheses of
Theorem~\ref{thm:geometric-closure}. That theorem develops the graph,
fills the inner disk, attaches the marked cap, and yields the family
\[
 F_\tau=R_{-2\pi\tau}\circ F\circ\varrho_{2\pi\tau}
\]
on the disk, with the inclusion on the cap. On the collars the two
rotations cancel pointwise; the marked cap makes two of the closed
images noncongruent; and the same theorem gives embeddedness and the
constancy of the area, of the enclosed volume and of the total
oriented mean curvature.
\end{proof}
